\documentclass[12pt,a4paper,onecolumn]{article}
\title{Digest of The New York Times}
\author{The New York Times}
% \definecolor{dark-green}{rgb}{0,0.25,0}
% \usepackage[utf8]{inputenc}

\usepackage{fontspec,xunicode,xltxtra}
\setmainfont[Ligatures=Required,Mapping=tex-text]{Georgia}
% \setmainfont[Ligatures=Required,Mapping=tex-text]{Palatino Linotype}
\setsansfont[Ligatures=Required,Mapping=tex-text]{Myriad Pro}
\setmonofont[Ligatures=Required,Mapping=tex-text]{Courier New}
\usepackage{xeCJK}
\setCJKmainfont{宋体}
\setCJKmonofont{黑体}

% \usepackage{indentfirst}


\usepackage[dvips,dvipsnames,svgnames]{xcolor}
\definecolor{light-gray}{gray}{0.95}

%生成PDF的链接
\usepackage[dvipdfmx]{hyperref}
\hypersetup{
    bookmarks=true,         % show bookmarks bar?
    bookmarksopen=true,
    pdfpagemode=UseNone,    % options: UseNode, UseThumbs, UseOutlines, FullScreen
    pdfstartview=FitH,
    pdfborder=1,
    pdfhighlight=/P,
    pdfauthor={wuxch},
    unicode=true,           % non-Latin characters in Acrobat’s bookmarks
    colorlinks,             % false: boxed links; true: colored links
    linkcolor=blue,         % color of internal links
    citecolor=blue,        % color of links to bibliography
    filecolor=magenta,      % color of file links
    urlcolor=cyan           % color of external links
}
% 使用目录
\usepackage{makeidx}
% 使用图形
\usepackage{graphicx}


\usepackage{titletoc}
\dottedcontents{section}[1.5em]{}{2.3em}{1pc}

% % 页眉页脚的设定
% \usepackage{fancyhdr}
% \pagestyle{fancy}
% \lhead{\includegraphics[height=0.2in]{zte_logo}}

% \usepackage{footmisc}
\usepackage{pfnote}
\usepackage{fnpos}

% % 嵌入的代码显示
% \usepackage{listings}
% \lstset{language=C++, breaklines, extendedchars=false}
% \lstset{keywordstyle=\color{blue},
%         commentstyle=\color{SeaGreen},
%         stringstyle=\ttfamily,
%         showstringspaces=false,
%         backgroundcolor=\color{light-gray}}
% \lstset{basicstyle=\small}

\usepackage{booktabs}

\usepackage{tipa}
\usepackage{phonetic}
\usepackage[gen]{eurosym}

% \usepackage{CJK}
\usepackage{makeidx}
% \usepackage{glossary}

\usepackage{fancyhdr}

\usepackage{lastpage}
\cfoot{\textsf{\thepage\ of \pageref*{LastPage}}}

\pagestyle{fancy}
\fancyhead{}
\renewcommand{\headrulewidth}{0pt}
\pagenumbering{arabic}

\usepackage{graphicx}
\usepackage{picinpar}

\setlength{\parindent}{0in}
\setlength{\parskip}{2.5ex plus4mm minus3mm}
% \setlength{\parskip}{.3\baselineskip}

\setcounter{section}{0}

% bibtex
\usepackage[super]{natbib}
% \usepackage[globalcitecopy,labelstoglobalaux]{bibunits}
% \makeatletter
% \renewcommand{\@bibunitname}{\jobname.\the\@bibunitauxcnt}
% \makeatother

% \renewcommand{\refname}{缩略语}
% \renewcommand{\refname}{}

\makeatletter                   % Reference list option change
\renewcommand\@biblabel[1]{#1:  } % from [1] to 1: 
\makeatother                    %

\newcommand{\citenumfont}[1]{\textsf{#1}}

% \usepackage{tipa}
% \usepackage{phonetic}

\newcommand{\phonetic}[1]{{\fontspec[Colour=CC00CC]{Microsoft Sans Serif} #1}}


%%% Local Variables:
%%% TeX-master: "medical_trivia"
%%% End:

\makeindex
\begin{document}
\date{}
\thispagestyle{empty}
\renewcommand\contentsname{\textsf{Digest of The New York Times}}
\textsf{\tableofcontents}
\pagebreak
\setcounter{page}{1}
\section{Mining Giant Scraps China Deal}

SHANGHAI -- The Chinese government's largest investment in a Western company, a proposed \$19.5
billion stake in the Australian-British mining giant Rio Tinto, collapsed Friday, dealing a blow to
China's global corporate ambitions and its efforts to gain clout in the natural resources market.

The board of Rio Tinto announced the decision after meeting in London on Thursday, saying the
company had ended the deal struck in February to sell the stake to the Chinese state-owned Aluminum
Corp.~of China, also known as Chinalco.

The board said in a statement Friday that it would instead raise about \$20 billion by issuing new
stock and forming a joint venture with its longtime rival, the Anglo-Australian mining giant BHP
Billiton.

The China deal, which would effectively have expanded Chinalco's stake of 9.3 percent in Rio to 18.5
percent, had drawn stiff political opposition in Australia, where mineral riches have fueled
prosperity and where some have begun to fear China's power in the region. Institutional investors
had also expressed skepticism about the agreement, with some calling it a sweetheart deal for the
Chinese government. Chinalco said it regretted the decision and had worked hard to try to revise the
deal to reflect changed market conditions, as well as the response from shareholders and regulators.

``As a result, we are very disappointed with this outcome,'' Chinalco's president, Xiong Weiping,
said in a statement.

In its statement, Rio Tinto said that it had agreed to raise \$5.8 billion through a venture with
BHP Billiton and that it planned to raise \$15.2 billion more through a stock offering. Rio agreed
to pay Chinalco a \$195 million breakup fee.

In Australia and elsewhere, the Rio-Chinalco deal has been viewed as a corporate landmark, a crucial
test of whether the West's shareholder-driven capitalism could mesh with a Chinese model that is
ultimately under state control.

The collapse of the deal will rank alongside the failure by a Chinese company in an attempt to take
over Unocal, the U.S.~oil company, in 2005. It will provide new evidence that Chinese attempts to
extend its global corporate reach are often limited by nationalistic reactions abroad.

The Unocal deal collapsed amid furious debate in the United States about the wisdom of selling
strategic assets to China and over China's human rights record.

Until recently, Rio and Chinalco had strongly defended the investment, which would have made the
Chinese government Rio's single biggest shareholder by far and given it an advisory role in the
company's operations. Chinalco would also have gained substantial stakes in individual mines in
several countries.

In return, Rio would have received a sorely needed injection of cash to retire a portion of its \$39
billion in debt, amassed mostly in an ill-timed 2007 purchase of the Canadian aluminum maker Alcan.
Rio would have formed a strategic partnership with Chinalco that company officials say would have
given it greater access to Chinese customers and opportunities to search for ores in China.

The deal was made last year at the depths of the collapse in global commodities and stock markets.
The share price for Rio had been severely depressed, not only by falling commodity prices but also
by a costly defense against a proposed and potentially lucrative takeover by BHP Billiton.

But the tentative agreement came under increasing fire as the global downturn waned, raising metals
prices and Rio's stock price considerably from their earlier lows. Last February, shares in Rio
Tinto were trading at about \textsterling18 in London; on Friday, they were at \textsterling30.94,
or \$49.60, in afternoon trading.

Institutional shareholders complained that Rio should raise money through a stock issue instead of
selling a dominant stake to the Chinese government, which they argued had an interest in keeping the
price of ores as low as possible.

Chinalco insisted that it was independent of the Chinese government, and Australian regulators ruled
that the Chinese would not be able to affect the price of iron ore, one of Rio's main products, even
with the proposed stake in the company.

But those declarations failed to sway some critics, who argued that Chinalco's proposed investment
followed the government's declared strategy of diversifying its global investments in resources as a
way of hedging against price increases. The Chinese are such large consumers of some metals, like
iron ore, that they already have some say over prices.

Opposition politicians in Australia savaged the deal, calling it evidence that Prime Minister Kevin
Rudd, a Mandarin-speaking former envoy to Beijing, was too cozy with the Chinese government. A
television commercial nationally broadcast in the past week showed scenes of the Chinese Army's 1989
assault on pro-democracy demonstrators in Tiananmen Square and charged that Mr.~Rudd ``bears gifts
to the Chinese military regime by allowing control of strategic mineral resources in Australia.''

A Nationals senator, Barnaby Joyce, one of the Australian politicians who opposed the Rio-Chinalco
deal, said the falling out was ``great for the Australian people.''

``We do not have the complications of the Communist Peoples' Republic of China's government owning
the wealth of Australia, in the ground in Australia,'' Mr.~Joyce told the Australian Broadcasting
Corp.~on Friday.

Mr.~Rudd said that the country was still open to Chinese investment, but that future bids would be
weighed against Australian national interests. He said that politics were not responsible for the
collapse of the deal.

As the debate over the Rio investment grew louder, Chinese officials complained about what they
called protectionism and nationalism. The Australian Foreign Investment Review Board, which
scrutinizes deals that involve big stakes in Australian companies, was expected to rule on the
Chinalco investment this month or early next month.

The proposed investment was seen as another sign of China's growing global clout and its need for
the ingredients that help produce the steel used in the country's fast-paced building boom. Chinese
companies have been on the prowl around the world over the past few years, securing supply contracts
and buying corporate stakes that secure long-term access to oil, iron ore, rare earth metals and
other natural resources.

The Chinese have also made major investments in recent years in real estate, agriculture and
international finance. Just in the past week, the Chinese sovereign wealth fund agreed to spend more
than \$1 billion to increase its already significant stake in Morgan Stanley, the Wall Street
investment bank.

Under the Rio deal, Chinalco had offered \$7.2 billion to purchase convertible bonds in Rio Tinto
and \$12.3 billion in cash for stakes in Rio's aluminum, iron ore and copper assets in the United States, Australia and Chile. The bonds could have been converted to stock.

\section{2 U.S. Journalists on Trial in N. Korea}

SEOUL -- Two American journalists were scheduled to go on trial on Thursday in a case that is widely
seen as a test of how far the isolated Communist state was willing to take its confrontational
stance toward the United States.

The two Americans, Laura Ling and Euna Lee, have been held in North Korea for two and one-half
months on charges of illegal entry and ``hostile'' acts. If they are convicted on such charges under
the Communist state's criminal code, the two journalists could be sentenced to five years of hard
labor or more at the country's infamous prison camps.

Secretary of State Hillary Rodham Clinton has called the charges ``baseless.'' Washington reiterated
on Wednesday that the North should forgo the legal proceedings and release the two women.

In a brief dispatch early Thursday, the North's official news organ, the Korean Central News Agency,
said that the trial would begin at 3 p.m. (2 a.m.~Eastern time on Thursday) in Pyongyang Central
Court, the highest court of North Korea. But by Thursday evening, there was no word on whether the
trial had begun.

Ms.~Ling and Ms.~Lee were on a reporting assignment from Current TV, a San Francisco-based media
company co-founded by Al Gore, the former vice president, when they were detained by North Korean
soldiers on the border between North Korea and China on March 17. The reporters were working on a
report about North Korean refugees -- women and children -- who had fled their impoverished country
in hopes of finding food in China. The circumstances surrounding the detention of the two
journalists remain unclear.

Analysts said they became a pawn in a rapidly deteriorating confrontation between the United States
and North Korea -- a potential bargaining chip for the Pyongyang regime and a handicap for
Washington in its efforts to pressure the regime over its recent missile and nuclear tests.

The reporters' families pleaded for clemency.

``All we can do is hope the North Korean government will show leniency,'' Lisa Ling, Ms.~Ling's
older sister who is herself a television journalist, in an emotional plea at a California vigil
Wednesday night, The Associated Press reported. ``If at any point they committed a transgression,
then our families are deeply, deeply sorry. We know the girls are sorry as well.''

Lisa Ling said that North Korean officials permitted her sister to call her last Tuesday night,
roughly two days after its nuclear test.

``She said that the only way she may be able to get released is if our two countries communicate,''
Lisa Ling told Larry King of CNN on Monday, describing the four-minute telephone conversation she
had with her sister. ``She sounds absolutely terrified.''

Defying not only its traditional foes -- the United States, Japan and South Korea -- but also its longtime ideological allies, China and Russia, North Korea launched an intermediate-range rocket on April 5 and conducted an underground nuclear test on May 25.

Washington is leading an effort at the United Nations Security Council to impose tough sanctions on the North, which the Pyongyang government said it would consider an act of war.

\section{Tiananmen Square Scars Soldier Turned Artist}

BEIJING -- Soaked in sweat, his heart racing, Chen Guang descended the steps of China's Great Hall
of the People and aimed his automatic rifle at the sea of student protesters occupying Tiananmen
Square. A 17-year-old soldier from the countryside, Mr.~Chen and his comrades had just been given
chilling orders: to clear the symbolic heart of the nation, even if it meant spilling blood.

``We were assured there would be no legal consequences if we opened fire,'' Mr.~Chen recalled in an
interview on Tuesday. ``My only hope was that the students would not put up a fight.''

Twenty years after Chinese troops shot their way into the center of Beijing, killing hundreds of
people and wounding many more, Mr.~Chen provided a rare window into the military crackdown that
re-established the Communist Party's supremacy after six weeks of mass unrest and then, for most
Chinese, disappeared in an official whitewash.

Speaking publicly for the first time -- and defying security officials who have told him to keep
silent -- he explained how soldiers from the 65th Group Army dressed in civilian clothes on June 3
and stealthily made their way to the Great Hall on Tiananmen Square's western edge. At midnight,
with clips of ammunition slung across their chests, they faced off against demonstrators, the air
filled with the singing of students and the sound of gunfire.

``I can assure you I didn't shoot anyone,'' he said.

Now an artist and a bit of a provocateur living on the outskirts of Beijing, Mr.~Chen said he spent
the next 20 years suppressing memories of that day. But last year he began working on a series of
paintings based on hundreds of photographs, taken at his unit's request while he was on the square.
They include gauzy images of protesters commandeering a public bus, exuberant students parading with
pro-democracy banners and soldiers feeding their abandoned encampments into bonfires.

``For 20 years I tried to bury this episode, but the older you get the more these things float to
the surface,'' he said, chain-smoking in his apartment. ``I think it's time for my experiences, my
truth, to be shared with the rest of the world.''

But by publicizing his experiences through his art, Mr.~Chen risks provoking the authorities, who
are eager to suppress discussion of the episode and excise June 4 from public memory. In recent
weeks, as the anniversary of the crackdown approached, the police have harassed or detained
dissidents who they feared might draw attention to June 4. Last spring, Zhang Shijun, a former
soldier from north China, was arrested after telling The Associated Press that he regretted his role
in crushing the pro-democracy protests.

Last summer, after local galleries refused to show his paintings, Mr.~Chen posted them on the
Internet. Within hours, however, they had been taken down.

A slightly built man who talks softly and without emotion, Mr.~Chen says he is not worried about the
consequences of speaking out, even if he has received warnings to keep his paintings to himself.

``I'm not doing anything wrong,'' he said. ``I'm just talking about my experiences.''

Raised in rural Henan Province, the son of a factory worker, he dropped out of high school at 15
because, he said, he was a poor student. He wanted to be an artist, but everyone told him that was
no way to make a living. ``The pressure from my family was intense so I decided to join the army,''
he said. Because enlistees had to be at least 18, he lied about his age.

Less than a year later, in mid-April, Beijing was convulsed by protests touched off by the death of
Hu Yaobang, the Communist Party chief who had been forced to resign to take responsibility for what
some rival leaders viewed as reckless economic and political reforms.

Isolated in their barracks three hours north of the city, Mr.~Chen said he and his fellow soldiers
understood little about protests. They knew only what military officers told them: ``that bad people
were trying to destroy the nation that was established with the death of martyrs,'' he recalled.

On May 19, they were given orders to enter the city. But their path was blocked by throngs of
students and average Beijing residents supporting the demonstrators. For two days the troops were
lectured to and fed by strangers while the nation's military leaders debated what to do.

On the third day, his unit withdrew, but Mr.~Chen said the episode left him confused. ``We were told
they were bad people but the students seemed so honest and earnest,'' he said.

After nearly two weeks isolated in their barracks, the soldiers were given civilian clothes and told
to make their way to the Great Hall in groups of two or three. Mr.~Chen said his assignment was far
more unnerving. He said he was the only passenger in a double-length bus with its seats removed and
its interior filled to the windowsills with guns and ammunition.

Unfed and terrified, the soldiers, most of them teenagers, waited inside the Great Hall while
military commanders, perched at a second-floor window, strategized the assault. Around midnight,
power to the square was cut and the soldiers eased their way down the broad steps to the street. To
frighten the students into leaving, he said, the men were told to fire into the air. The tactic had
the desired effect.

By 2 a.m., tens of thousands of students were weeping and singing the Internationale as they filed
out of the square. Not long afterward, armored vehicles rolled in. One went to work on the Goddess
of Democracy, a papier-m\^{a}ch\'e statue that art students had built just days earlier. ``It took them
three rams before it fell to the ground,'' Mr.~Chen said.

Most of the deaths in the crackdown, according to multiple accounts of the incident, occurred in the
streets leading toward the square, not in the square itself.

Less than a year after the suppression, Mr.~Chen enrolled in the military's art school, then
transferred to the Chinese Academy of Fine Art. In 1995, he left the army.

In those early years, Mr.~Chen was drawn to photography and performance art, creating work that was
lurid and provocative. He spent months filming prostitutes and took photographs of himself
copulating on the Great Wall. He also produced a series of sexually explicit photos of himself
posing with an elderly intellectual man who had been persecuted by the Communist Party. ``I wanted
to portray myself having a visceral connection to someone who had experienced China's tumultuous
history,'' he said.

Although none of his early work refers directly to Tiananmen Square, he said most of it had been
influenced by the trauma there. ``Even if a connection is hard to see, everything I do is touched by
that experience,'' he said. Mr.~Chen said he saw soldiers bloodied by rocks and a protestor having
his head rifle butted by soldiers. But the image that haunts him most is rather mundane. As he was
cleaning up the square that morning, he spotted a luxuriant ponytail amid the detritus of crushed
bicycles and tangled blankets.

The clump of hair, held by a purple band, had been crudely shorn, perhaps as an act of protest but
possibly the result of something more sinister. ``It was a startling image,'' he said. ``I can't
stop thinking about that hair and why it had been cut off.''

In recent months, he has produced a score of self-portraits. In each, his neck, shoulder and chest
are littered with scraps of hair. He cuts his own hair only every year or two and then stores the
clippings in his apartment. So far he has filled the equivalent of two dozen coffee cans, raw
material for a future project.

He said he did his most intense work every June, around the same time that he was hit with wrenching
stomach pain. It is the same twisting of the gut that he first experienced on the square, he said.

It was around this time last year that Mr.~Chen decided to revisit the cache of photographs he had
taken in 1989. Just before the assault on the square began, Mr.~Chen's commander handed him a camera
and 20 rolls of film and told him to wander freely. When it came time to hand back the film, he hid
three rolls in his army-issued satchel.

He said the photographs inspired him to take on a subject that few in China care, or dare, to touch.
His paintings are artistic depictions of history, he insisted, not expressions of right or wrong.
The images are largely dispassionate, although Mr.~Chen has rendered them in a washed-out,
melancholy blue.

``I have no regrets about what I did,'' he said. ``But I feel that this tragedy could have been
avoided. Maybe if we start talking about this event, we can prevent it from happening again.''

\section{After Tiananmen and Prison, a Comfortable but Uneasy Life in the New China}

BEIJING -- When Liu Suli was released from a Beijing prison in 1991, having served 20 months for his
role in the pro-democracy Tiananmen Square protests of 1989, he pledged to abide by what he called
``the three noes.'' He would grant no interviews about the protests. He would write no articles. He
would accept no donations from sympathizers.

On that warm June 4 morning, Mr.~Liu carried off the corpses of protesters, their blood spattering
his shoes. Their sacrifice was too painful to face, much less reconcile.

But this week, as the 20th anniversary of the democracy movement's violent end drew near, Mr.~Liu
relented. Seated in his Beijing bookstore, interrupted only by countless drags on unfiltered
cigarettes and sips of chilled tea, he tried to reconcile Tiananmen then, and China now. He spoke
for three hours.

``If I don't talk about it any more, I might forget,'' he said. ``I might not know how to say it.''

Thursday marks 20 years since hundreds of students, workers and average citizens died in an army
assault on the Tiananmen Square democracy movement in the center of Beijing. Inside China, the day
is a nonevent; even oblique references to June 4 are taboo.

For the outside world, Chinese officials this year offer a tenable excuse: the protests, they say,
threatened China's stability. With stability restored, China produced an economic miracle that has
lifted scores of millions from destitution.

For veterans of the June 4 movement, however, reconciliation is seldom so tidy or superficial.

Some, Mr.~Liu says with evident disgust, have put the past aside and embraced the China where Adam
Smith, not Marx, is chief ideologist. Others have become permanent dissidents.

And some dwell in an uneasy middle world, part of both Tiananmen and the new China but not wholly in
either. They are prosperous lawyers who defend dissidents knowing that it could cost them their
jobs, and comfortable academics whose names on a human-rights petition could doom them to demotion
or exile.

And they are intellectuals like Liu Suli, who weathered the Tiananmen assault and prison to take a
place he never sought in China's new middle class. Over the 18 years since his release, he has built
a popular bookstore and cafe near Beijing's top academic institution, Peking University, that has
become a haven for intellectuals and expatriates. Now 49, he does not seem a man who has suffered
these 20 years of stability.

He suffers nonetheless. He is appalled by the notion that China's success was erected on the bodies
of Tiananmen demonstrators. He struggles with the inability of Tiananmen's ideals to gain more than
a toehold in his homeland. And he despairs at the oft-expressed thought that his fellow Chinese have
made a devil's bargain, trading the freedom that he and his fellow protesters sought for a chance at
a car and a bigger apartment.

``You can raise pigs to be very strong and very fat,'' he said. ``But a pig is still a pig. And a
pig has no rights.''

Mr.~Liu was 29, with a wife, a 2-year-old son and a lecturer's job at the Chinese University for
Political Science and Law when Hu Yaobang, the Communist leader who was pushed from power for
advocating liberal reforms, died in April 1989. Within a week, 100,000 demonstrators filled
Tiananmen Square for his funeral. One of Mr.~Liu's students, Zhou Yongjun, became the new democracy
movement's first leader.

Mr.~Liu was there too, but as an observer. Camera in hand, he shot dozens of rolls of film
documenting the events, and he began to see a timeline of sorts leading to Tiananmen, from protests
in 1986 and 1979.

``The political demands the students were making were actually an accumulation of 40 years of
Communist rule in China,'' he said. ``The na\"ivet\'e and idealism of the students wasn't just the
students; it was the na\"ivet\'e and idealism of the people.''

By the time the students began a hunger strike on May 13, the government had already ruled their
protests illegal. China's rulers, Mr.~Liu said, had left themselves no option but violence should
the students refuse to back down.

Eventually, he shed his outsider's position and joined the demonstrators, quickly rising to become
the operations director and press spokesman for the Capital Alliance to Protect Law and Patriotism,
an ad hoc group that coordinated the protests. He escorted Communist Party officials into the crowd
to talk to demonstrators, briefed Western journalists and escorted student leaders to meetings with
Chinese Army officers.

On June 4 soldiers approached the square's monument, where Mr.~Liu was stationed. One put a gun to
his head. An officer approached from behind, warning him in a low voice to flee.

They retreated through a chaotic field of overturned buses and burning military vehicles, returning
to a university campus where the bodies of seven protesters had been laid on the lawn. He helped
move them out of the sun to slow the decomposition process. Mr.~Liu and two fellow protesters,
wanted men, fled Beijing at the urging of protest leaders.

But Mr.~Liu soon returned. ``If everybody in the movement escaped,'' he said, ``there would be
nobody left to continue the movement.'' He spent a year and eight months in a Beijing prison with 13
other men, so crowded that not all could lie down at night to sleep. His wife, from whom he has been
since divorced, left with his son for the United States.

Mr.~Liu's release in February 1991, even before his case reached trial, was a welcome shock. He was
fired from his university position and expelled from the Communist Party, but that suited him. ``I
told them that anything that has to do with the party has nothing to do with me,'' he joked.

He says he is still watched constantly by security agents. ``The government really hates us,'' he
said. ``It pains them that we are here. We do a lot of things they don't like. But the price of
destroying us is something they can't afford.''

He says he is now devoting his energies to deeper discussions on how China can rejoin the modern
democratic world. For someone who saw his hopes crushed at Tiananmen, it is a tall order, and he
knows it.

``An idealist in China is a word that stinks up the streets,'' Mr.~Liu said. `` 'Idealist' equals an
insane person, or a stupid person, or an idiot. It's best that you don't talk about idealists in
China.''

So is he an idealist?

``Of course I am,'' Mr.~Liu said.

\section{North Korean Leader Is Said to Pick a Son as Heir}

WASHINGTON -- American and South Korean officials say that Kim Jong-il, the North Korean dictator
recovering from a stroke, appears to have designated his youngest son as his successor. Mr.~Kim may
have ordered the country's second nuclear test last week in hopes of leaving his son in control of a
country that has cemented its status as a nuclear-weapons state.

Little is known about the youngest son, Kim Jong-un, beyond reports that he was secretly schooled in
Switzerland under an assumed name, posing as the son of a driver in the nearby North Korean Embassy.
He skied Switzerland's famous slopes, and was a fan of Michael Jordan.

But his path to power is hardly assured: some intelligence officials believe that everyone from the
North Korean military to Kim Jong-il's eldest son may be plotting behind the scenes to derail the
succession plans, and North Korea's last ally, China, is reportedly deeply uncomfortable with the
thought of a third-generation family dynasty, unique among Communist nations.

As always in watching a state virtually sealed off from the outside world, analysts acknowledge that
they are extrapolating from indicators, rather than hard evidence. Even Kim Jong-un's age is
uncertain; he is believed to be in his mid-20s.

In recent weeks, North Korean diplomats abroad have been told to begin to pay homage to Kim Jong-un
and some schoolchildren have reportedly been including his name in their songs. His rise comes at
the expense of his 38-year-old brother, Kim Jong-nam, best known for the moment when he was caught
slipping into Japan on a false passport, on his way to Tokyo Disneyland.

In preliminary assessments of the May 25 underground nuclear test, top officials of the Obama
administration say they believe it was linked to the power jockeying in Pyongyang more than any
attempt to force President Obama into more negotiations in which North Korea tries to trade away its
nuclear ability for American energy and security concessions.

``There was a sense that every North Korean escalation was intended as a bargaining chip,'' said one
senior administration official involved in the assessments. ``Now there's an alternative view taking
hold: that Kim Jong-il wants to force the world to acknowledge it as a nuclear power before he dies.
And we're not going to do that.''

Mr.~Obama has sent James B.~Steinberg, the deputy secretary of state, and other senior officials to
Japan, South Korea and China this week to attempt a unified response to the test, focusing on
inspecting all sea and air traffic moving in and out of North Korea that could be carrying nuclear
or missile technology, according to an administration official.

But it is a risky move, one that could leave Mr.~Obama facing a possible military confrontation with
the North as he increases troop levels in Afghanistan.

The most delicate issue is whether China will back the inspection of some North Korean vessels,
which the North has already warned it would consider an act of war. The inspections were authorized,
but not enforced, under a United Nations resolution passed after the North's first nuclear test, in
2006.

A senior administration official said Tuesday that Chinese officials feared that if North Korea
gained the ability to fit its nuclear weapons atop its long-range missiles, the United States would
increase its military presence in the Pacific and Japan could reconsider its ban on its own nuclear
weapons. But if the Chinese press too hard, the official said, ``they risk unintentionally causing
collapse in North Korea and instability'' on their own border. The official spoke at a forum on
Tuesday under rules that he not be quoted by name.

While Obama administration officials say they want to draw North Korea back into disarmament talks
begun by President George W.~Bush, they are not interested in negotiating yet another deal to
disable the main nuclear facility at Yongbyon, where North Korea produces its bomb-grade plutonium.
``The real challenge is to avoid a repetition of the past,'' one senior administration official
said.

Obama administration officials acknowledge that negotiations with the North may be all but
impossible at a moment when it is unclear who is running the country and when all players in a
succession struggle will avoid any perception of concessions to the United States.

There is no indication yet that the heir apparent has been involved in decisions about the nuclear
program.

The current leader, Kim Jong-il, has three known sons. The eldest, Kim Jong-nam, was once considered
the leading candidate to succeed his father, until the Disneyland episode added to rumors that his
judgment was less than reliable. Kim Jong-nam is widely reported to have voracious appetites for
alcohol and women, and his father apparently grew concerned that North Korea's generals would never
accept him, according to a former American intelligence official.

The North Korean leader's middle son, Kim Jong-chol, 28, was another possibility, but Kenji
Fujimoto, who once served as Kim Jong-il's sushi chef, wrote in a memoir that Kim Jong-il dismissed
that son as ``girlish,'' suggesting that he would not stand up to the West.

By default, that left Kim Jong-un. On Tuesday, South Korean lawmakers said they had been briefed by
the country's intelligence agency, and told he was the heir apparent. The intelligence agency
intercepted messages to North Korean overseas missions a few days after the May 25 nuclear test,
according to reports in Seoul.

``Our intelligence service has been following the matter for some time,'' said Song Young-gil, an
opposition lawmaker briefed by the intelligence agency. ``They said that this message instructed the
diplomats to pledge their allegiance to Kim Jong-un.''

Another lawmaker, Moon Kook-hyun, said he could not comment on a secret briefing but agreed that Kim
Jong-un had been designated the successor. The intelligence agency declined to confirm the reports.

Inside the North, the subject is only whispered about. ``I never thought that Kim Jong-il was human
and thus mortal,'' said Oh Yeon-jong, a defector who arrived in Seoul in 2004. ``We didn't know,
didn't talk about how many children he had, how many wives he had. I heard about them only when I
arrived here.''

It also is not clear if a society that reveres seniority would accept such a young leader, and
American officials are waiting for the next steps.

One test could be the response of Gen. O Kuk-ryol, the National Defense Commission's vice chairman.
Intelligence officials say they believe he would need to give his blessing to the transfer of power
in Pyongyang. The general has taken greater control in recent years over the regime's military and
security policy. Analysts are also watching Jang Seong-taek, Kim Jong-il's brother-in-law, who is
believed to run many day-to-day state affairs on behalf of the ailing leader.

In any case, the outside world will be transfixed. At the C.I.A., the youngest son's picture used to
be posted in the Asia division, and analysts gave him the moniker of ``the Cute Leader'' -- a play
on his father's status as ``Dear Leader'' and his grandfather's as ``Great Leader.''

\section{Apps Boom as Companies Seek a Place on Your Phone}

Developers of programs for the iPhone have already managed to make a decent living selling hundreds
of thousands of copies of games from their living rooms or garages.

But now, a new way to profit from writing software for the iPhone is emerging: Sell the apps, then
sell your company.

With the number of downloads through Apple's App Store topping one billion and more than 40 million
iPhones and iPod Touches sold since 2007, an increasing number of companies are seeing the mobile
industry as a source of sustained revenue. Recently, IAC/InterActiveCorp, the Internet media
conglomerate founded by Barry Diller, and Amazon.com, have bought app developers. Smaller companies
have begun to assemble properties.

Since Apple showed that new apps sell phones, the market for apps is expanding quickly. Palm,
Research in Motion, Nokia and Microsoft are all building app stores to work with phones running
their operating systems. Apps can also be built for phones running Google's Android software.

Most of the action is still in iPhone apps, which is what makes Apple's Worldwide Developer's
Conference this week in San Francisco of interest to developers and potential investors.

Developers will be showing new products running on Apple's latest software, which allows users to
buy subscriptions to applications and easily buy add-ons like access to higher game levels or
additional city guides. The potential for added revenue should increase interest from buyers looking
for acquisitions.

``There's going to be a lot more interest in iPhone applications after the upgrade,'' said Greg
Yardley, a co-founder of Pinch Media, a mobile analytics firm. ``We're going to see some really neat
business models emerge because of the new ability to sell virtual goods.''

The increased interest in app developers is being driven by companies seeking to build cellphone
apps for their products or services. They see it as a way to reach beyond the Web for consumers.
Though many apps are free, the willingness of people to pay 99 cents or more for one gives companies
hope that apps may be a more reliable source of revenue than Web sites.

``Companies are asking themselves, 'How can we get on the iPhone?' '' said Matt Murphy, a partner at
venture capital firm Kleiner Perkins Caufield \& Byers, which maintains a \$100 million fund devoted
solely to investing in start-ups creating apps for the iPhone. ``Instead of trying to organically
create their own property, they're looking at applications with traction and cherry-picking the ones
that seem like a good fit.'' (A polished, professional iPhone application can cost around \$50,000
to produce.)

With an instantaneous and established presence on the iPhone platform, he said, a company could tap
into a stable, loyal fan base. For a big company that is trying to go mobile, and quickly, ``those
few million users are almost more valuable than the property itself,'' he said.

That was the approach taken by IAC, which has more than 35 Internet-based companies, including
Ask.com, CollegeHumor and Evite. Last month, the company bought UrbanSpoon, a start-up based in
Seattle that recommends nearby restaurants, for an undisclosed sum. It is one of the App Store's
most popular products, having been downloaded close to five million times.

``The iPhone is a big part of our mobile strategy,'' said Leslie Cafferty, a spokeswoman for IAC.
``It takes a lot to invest in developing an application. It was much more appealing to pick one
up.''

CitySearch, another of IAC's properties, first worked with UrbanSpoon to syndicate advertisements
and reviews through the start-up's Web site. When the iPhone application made its debut last July,
it was ``the icing on the cake,'' said Ms.~Cafferty.

For the three creators of UrbanSpoon, it was ``partially opportunistic'' to be bought by a larger
company, said Ethan Lowry, one of the founders. ``It let us think on a grander scale of the services
we can offer.'' And it offered stability, he said. ``The security of a larger company has taken some
of the financial stress out of the situation,'' he said.

Amazon, with its Kindle electronic book reader, looked at phone apps as a way to expand the market
for the e-books it sells. In late April, Amazon bought Stanza, a software service that allows users
to browse and buy from a library of 100,000 books through a phone.

But it is not just larger companies beginning to see promise in the popularity of the phone and high
demand for its programs.

Tapulous, a start-up in Palo Alto, Calif., bought a game called Tap Tap Revenge from its developer,
Nate True, in July. The game, patterned after video games like Guitar Hero that challenge players to
keep rhythm with popular songs, has been downloaded by one out of every three digital shoppers in
the iTunes App Store, according to market research firm comScore.

``We're very excited to connect with people at W.W.D.C. who have interesting ideas and applications
in the music gaming space,'' said Andrew Lacy, the chief operating officer of Tapulous.

Other companies are making a business out of acquiring raw programs to redesign, polish and release
into the App Store. For example, Ngmoco, a video game start-up based in San Francisco devoted solely
to publishing games for the iPhone and iPod Touch, recently bought the quirky, colorful puzzle game
Rolando designed by a British developer, Simon Oliver.

``When we built our company, we decided the best model was to harness interesting developers and
bring their projects to market,'' said Neil Young, the chief executive at Ngmoco. Mr.~Young said the
company also develops games in-house, but is always reaching out to developers to see what they are
working on next.

So far, their approach appears to be working: many of their games rank highly in the App Store's
most-popular lists. The company has four more games to be released in the next two months and 15 new
games under development, including several sequels to Rolando.


\section{The Woman in Charge of Reality Shows at Spike TV}

LOS ANGELES -- Sharon Levy is the first to admit that she is not what anyone would call a girl's
girl -- just ask her, for example, about her television viewing habits.

``Grey's Anatomy?'' she repeats, incredulously. ``Kill me now.''

Which is exactly the reaction one might expect from the person in charge of developing reality
programming for Spike TV. That is the cable channel whose slogan, ``Get More Action,'' sums up the
sometimes sophomoric, always male-centered programming that has cemented Spike's reputation, fair or
not, as a sort of Maxim magazine for television. Ms.~Levy, 39, is out to refocus that, if not
completely change it. To do so, she has introduced a slate of reality shows this year that has
attracted new male viewers without depending so heavily on the stereotypes that have clung to the
cable channel since 2003, when it changed its name to Spike and began marketing itself as the
channel for men.

``She's got as much testosterone as any guy I know,'' said Doug Herzog, the president of MTV
Networks Entertainment Group, a part of Viacom and the parent of Spike. ``She was standing in the
hallway recently, talking about 'Wolverine' and 'Star Trek.' She's into video games and that geeky
guy stuff. At the same time, she will turn around and tell you how awesome her new handbag is.''
(She says she is particularly proud of one recent purchase, a limited edition Takashi Murakami/Louis
Vuitton.)

Among the new series she has brought to Spike this year is ``Deadliest Warrior,'' which enlists
computer modeling to determine which of two archetypal fighters would win a battle between, say, a
ninja and a Spartan or an Apache and a Roman gladiator.

Another, ``4th and Long,'' follows overlooked professional football hopefuls competing for a chance
to be invited to the Dallas Cowboys training camp. And ``Jesse James Is a Dead Man'' features
Mr.~James, the owner of West Coast Choppers and the former star of ``Monster Garage,'' as he
undertakes daredevil stunts.

The results have so far been noticeably good. ``Deadliest Warrior'' drew an estimated 1.7 million
viewers to its debut in April, 50 percent larger than the cable channel's average prime-time
audience and making the program at the time the highest-rated new cable series this year among young
men. (Only about 5 percent of the show's audience was women.)

``4th and Long,'' whose debut was scheduled opposite ``WWE Raw,'' the wrestling program that is one
of the highest-rated series on cable, nevertheless drew 1.2 million viewers. And the first episode
of ``Jesse James'' drew two million viewers last month, more than any other unscripted premiere in
the cable channel's history. Another 1.1 million tuned in for the rebroadcast the same night.

David Campanelli, vice president for national television at Horizon Media, said the recent additions
have done a better job of reaching Spike's intended male audience than the heavy slate of ``C.S.I.''
repeats that the channel had been scheduling. ``I wouldn't call any of it highbrow, but it's not
going after the lowest common denominator, either,'' he said.

Kevin Kay, the president of Spike, said he believes that the new series will be successful at
generating water-cooler talk, to spark and perhaps settle arguments about ancient warriors or
recount cool explosions.

That the water cooler should not be too near the human resources department, however. ``Manswers,''
one of Ms.~Levy's early creations, features breathless inquiries into topics like the best position
in which to have sex while scuba diving and whether breasts bounce differently on the moon.

While Spike in April ordered a third season of ``Manswers,'' more recent shows are moving further
away from the jiggle factor, Ms.~Levy promises, and more toward shows like ``DEA,'' which follows
the efforts of the federal Drug Enforcement Administration.

``We have been trying to figure out how to do a show with women in it that is sexy and not
misogynist,'' Ms.~Levy said. ``We have to crack that code.''

She believes that she is getting closer. A new series called ``Surviving Disaster,'' which provides
instruction on how to make it through the aftermath of a nuclear explosion or a hijacking, will draw
women even though it is not geared toward them, she predicted.

``This is going to be a very broad show,'' she said. ``We're still a network that targets men 18 to
34, but we can do a show that doesn't alienate women.'' The series is scheduled to make its debut
next year.

Long before she found herself as a woman at a men's cable channel, Ms.~Levy had grown accustomed to
assimilating. Raised in a kosher Jewish home in East Meadow, N.Y., in third grade she fatefully
traded lunch with a friend.

``She had a ham sandwich with mayo on it,'' Ms.~Levy recalls. ``I took one taste and thought, 'This
is amazing!' Then, at a sleepover party, I had real bacon. Suddenly I'm really into this pork thing.
Next thing you know, I fell off the religious wagon.'' Her husband proposed to her in front of their
Christmas tree. And both she and her toddler son love big trucks.

After completing a degree in European history at the State University of New York at Albany and
considering law school, she went to work in the publicity department of Comedy Central, a part of
Viacom, then moved to the Sci Fi Channel.

She was lured back to Comedy Central to help start a promising but vulgar new animated series titled
``South Park''; she also worked on ``The Man Show,'' which closed most episodes with well-endowed
women in tight T-shirts bouncing on a trampoline.

Soon, Stone Stanley Entertainment, the producer of ``The Man Show,'' hired Ms.~Levy as an executive
to develop new shows; a few years later, in 2005, she returned to the Viacom fold, this time at
Spike.

Along the way, Ms.~Levy struck up the relationships that she says have contributed to her view about
what kind of shows men like. Mr.~James, for example, cites Ms.~Levy as a friend to whom he had been
talking ``for a couple of years'' about a potential show.

``She wanted to do a 'Project Runway' with bikes,'' Mr.~James said. ``I kind of laughed at her. But
we kept working on it and the right thing came along.

``She's trying to survive in a man's world, so she's got it tough,'' Mr.~James said. ``She's like
that cool chick in college that you drink with and go out and party with. She has her girly moments,
but not too often.''

\section{For the Obamas, Date Night in Paris}

PARIS (AP) -- Ah, Paris. The city of love. And the city of this week's presidential date night.

A week after flying to New York for dinner and a Broadway show, President Barack Obama and first
lady, Michelle, dined at a cozy neighborhood bistro just a few blocks from the Eiffel Tower.

The president and first lady were in France to join their counterparts from France, Canada and
Britain to commemmorate the 65th anniversary of D-Day, the allied invasion of Normandy in World War
II that led to victory in Europe.

La Fontaine de Mars dates back to 1908 and specializes in rustic dishes from France's southwest
region of Bordeaux Perigord and the Basque. Foie gras, duck and cassoulet are on the menu, although
the White House has not said what the president and first lady chose.

The Obamas shared the restaurant with other diners, and other restaurants on the winding Rue St.
Dominique were filled. Police, some in riot gear, lined the street. Crowds pressed behind barriers
at the end of the street to glimpse the first couple, and about 100 people gathered there burst into
applause as the Obamas left the restaurant. Clusters of people at street corners held up cell phones
and cameras to snap pictures.

After dinner and a ride along the quai on the Left Bank of the Seine River, the Obamas returned to
the U.S.~ambassador's residence, where first daughters Sasha and Malia awaited them.

Earlier in the day, Sasha and Malia joined their parents on a visit to Notre Dame Cathedral, where a
children's choir sang and the president lit a candle. They climbed the stairs for a view of the city
from the roof of the 12th century Gothic church. The cathedral was closed to the public during the
first family's visit.

The first family's motorcade traveled to the cathedral along the Boulevard St. Germain, where
hundreds lined the street hoping to see them.

\section{What Would Mao Drive? A Little Red \ldots Hummer}

Even in a world mostly done being amazed by the ironies of globalization, last week managed to
produce something fresh and previously unfathomable: General Motors, newly bankrupt and struggling
to raise cash, agreed to sell its Hummer division to a company from China.

Yes, that Hummer, maker of the famously gas-guzzling behemoths whose menacing width and armor trace
their provenance to the American military, is now set to become the property of Sichuan Tengzhong
Heavy Industrial Machinery Company, in a land officially still called the People's Republic of
China.

It might seem incongruous, this plaything for the unabashed American road warrior shifting to a
country where the bicycle once ruled and collectivism was an organizing principle. (What next?
Harley-Davidson snapped up by the Vietnamese?) But that's just until you contemplate the realities
of modern China, and the nouveau riche in the growing suburbs, setting down lawn furniture inside
gated complexes of villas, shopping at big-box stores and driving luxury cars. China seems intent on
nurturing the very sorts of landscapes and consumer attitudes that produced the Hummer.

More than a merely economic event -- the latest sign of China's rise and American struggles -- the
deal is a cultural moment. It seems no accident that a Chinese company is taking possession of
Hummer. China has come to embrace many of the attributes and modes of consumption that Americans may
reflexively consider their own, complete with the sprawl and tangle of highways familiar to any
resident of Los Angeles or Atlanta.

As China has cast off its ideological past and aggressively modernized its cities, it might
reasonably have been expected to look to Europe or Japan for models of urban planning. Like Japan --
home to one of the most sophisticated rail networks on earth -- China is densely populated and
dependent on imported oil. As is true in Europe, China's major cities are surrounded by productive
agricultural lands, making tightly clustered growth seem prudent.

Instead, in a choice familiar to Americans, China has put the automobile at the center of
contemporary life. China has torn down older buildings in every major city to make way for more
vehicles. It has erected an impressive network of highways crisscrossing the vast country. Air
quality and energy efficiency have been outweighed by reverence for the car.

This has not happened randomly. In recent times, China's leaders have unleashed enormous quantities
of state finance to seed auto ventures in every province, spurring industries that have grown along
with the ubiquity of the car. Petrochemicals, steel, glass-making and rubber have all expanded to
feed auto-making. Tourism and retail shopping have increased as more Chinese take possession of
steering wheels.

Along the way, many Chinese aspirations have come to focus on car ownership. In a country where so
many people look back with bitterness on the regimented days of Maoism, and where public
transportation still involves packing into belching buses and gruesomely crowded trains, the car has
become a vessel for Chinese dreams.

``Why do you want a car?'' I asked a young professional couple shopping at a car lot outside Beijing
in 2002. The question elicited an irritated glare from the woman, as if I were condescending. ``Same
reason you want a car,'' she said. ``We want what you want.''

She did not mean merely the ability to go where she pleased, but also the geography the car enables
-- the villas with their backyards and modern conveniences; the superstores selling microwave-ready
food; the new golf courses.

Asked about the atrocious smog blanketing Beijing and the traffic jams that made driving there a
threat to mental health, she shrugged. Anyway, the car takes a growing slice of wealthy urbanites
into the countryside for the pure pleasure of it -- this, too, once unthinkable in a nation that
remembers how city dwellers were forcibly dispatched to the hinterlands to slop pigs during the
Cultural Revolution.

As driving has evolved from a mere way to get around to a mode of life, sales of passenger cars in
China have grown by 20 percent to 30 percent per year since 2005. In January, for the first time
ever, China's monthly vehicle sales exceeded those in the United States.

Still, it would be wrong, if tempting, to view China as some sort of time capsule of 1950s American
suburbia, that age before the vernacular came to include ``climate change'' and ``non-renewable
resource.'' China is moving to impose fuel efficiency standards tougher than those in the United
States while developing hybrid cars. Most car owners favor cheaper, smaller models.

But much as everywhere else, the auto has become a status symbol in China, and in ways that collide
with the country's history as a paragon of anti-imperialism: Cadillacs, that icon of American capitalist success, are now made on Chinese soil along with Audis and Buicks. In the new Chinese social divide, the middle class hews to economical vehicles with tiny engines, while the wealthy fuss with automatic climate control in their luxury sedans.

So it actually seems fitting that Hummer will be a Chinese brand. A vehicle that makes sense only as a pastime, with its miserable gas mileage, the Hummer has become an object increasingly shunned in the United States as a sign of wasteful decadence (not to mention something that most Americans can no longer afford). China, primed to consume and long since shorn of its collectively imposed thrift, will take a crack at extracting profit by selling the hulking beasts.


\section{iPhone Stars in Apple Show, Supported by Software}

SAN FRANCISCO -- Apple offered its devoted fans their periodic supply of catnip on Monday with a new
version of the iPhone, called the iPhone 3G S.

Apple announced some things that were characteristically Apple: more, better, faster features on its
smartphones and computers. It also announced something that was not: lower prices.

The company said the third model of its hit phone, which looks physically identical to the last
version but includes internal hardware and software improvements, will go on sale on June 19. Among
the changes, the iPhone 3G S will have a three-megapixel camera that also records video, an internal
digital compass and voice-control features that let owners use spoken commands to make calls and
play music.

The S stands for speed, and upgraded internal components will allow the device to run more than
twice as fast as the previous model, according to Philip W.~Schiller, Apple's senior vice president
for worldwide product marketing. Mr.~Schiller delivered much of the keynote presentation at Apple's
annual Worldwide Developers Conference, a meeting of Apple software makers that will last through
Friday.

Steven P.~Jobs, Apple's chief executive, was not present, as some commentators had expected.
Mr.~Jobs is on a six-month medical leave and is not expected to return until later this month.

The iPhone 3G S, which will be available in both black and white, is priced at \$199 for the
16-gigabyte version and \$299 for the 32-gigabyte version. A spokesman for AT\&T, Apple's exclusive
iPhone partner in the United States, said that for now, it would not change the monthly prices on
its iPhone voice and data plans.

Apple also announced it would sell the previous eight-gigabyte iPhone 3G for a price of \$99,
reduced from \$199, on top of data and voice plans. That could make the iPhone accessible to
cost-conscious consumers and undercut competitors like Palm. Palm's new smartphone, the Pre, which
went on sale Saturday, costs \$199 after a mail-in rebate. Sprint, which is selling the Pre
exclusively, offers a two-year service contract that is cheaper than the one AT\&T offers for the
iPhone.

``This puts a lot of pressure on the likes of RIM and Palm,'' said Charles Wolf, an analyst at
Needham \& Company, referring to the BlackBerry maker Research in Motion. ``It sorts of upsets
Palm's apple cart two days after the launch of the Pre.''

Apple also used the industry gathering to announce new features and price cuts on its MacBook Pro
line of laptops. All new MacBook Pros will feature up to seven hours of battery life -- a 40 percent
increase from the last version, because of the new lithium-ion batteries Apple is using. The new
batteries are sealed into the laptops, as opposed to the removable batteries that have been in
previous MacBooks. Apple is saying the power units can be recharged up to 1,000 times -- up to three
times as much as other batteries. New MacBook Pros will also feature an SD memory-card slot for
simpler photo transfers.

The new, lower prices suggest the company is taking aim at recession-wary consumers and is willing
to sacrifice at least some profit to increase its share of the market for personal computers.

But Apple did not budge on prices significantly. For example, the new 15-inch MacBook Pro now costs
\$1,699, \$300 less than the previous version.

``I think there was a recognition that their prices are becoming increasingly disconnected from the
market,'' said Toni Sacconaghi Jr., an analyst at Sanford C.~Bernstein \& Company. ``But this was
not a bold step into new price points.''

AT\&T said iPhone owners with existing contracts would have to pay more to upgrade to a new phone:
\$499 for the eight-gigabyte, \$599 for the 16-gigabyte and \$699 for the 32-gigabyte. However,
current customers who are eligible for upgrades and those whose contracts have expired would pay the
same price as a new customer, said Mark Siegel, executive director of media relations for AT\&T
mobility and consumer markets. The original iPhone was sold two years ago.

TomTom, the maker of GPS systems for cars or personal use, also made a presentation at the
conference, in which it said it was introducing a portable navigation application and car kit to
provide drivers with turn-by-turn directions. TomTom executives were acknowledging that the iPhone
and other smartphones were killing its market for stand-alone GPS hardware and that the company was
entering the accessories and applications business.

The new software for the iPhone also pushed the device more clearly into the portable game console
arena occupied by Nintendo's DS and Sony's PSP. In what looks like another play for the iPhone to
resemble a handheld gaming device, the new peer-to-peer connection feature will locate nearby users
using Bluetooth and connect the two using Wi-Fi, allowing them to compete in various games.

Much of Apple's presentation was devoted to software, both for the Macintosh computer and the
iPhone. The company said a new version of its operating system for the Macintosh, called Snow
Leopard, would be available in September. It will cost \$29 for current Mac owners to upgrade and
the operating system will be shipped with all new Macs.

Snow Leopard includes new features that include enhanced global search that can now read URLs and
text in Web pages; better compatibility with Microsoft Exchange, the dominant e-mail system used in
businesses; and faster computing over all.

Apple also said the new version of the iPhone's operating system, iPhone OS 3.0, will be made
available free to current iPhone owners beginning June 17. The software will have more than 100 new
features, including capabilities for multimedia messaging; cut, copy and paste tools; and search
functions. It will also unlock features developers have eagerly awaited, like the ability to buy
virtual goods or new gaming levels within applications.

ScrollMotion, a start-up based in New York, unveiled an application at the convention called
Iceberg, a digital bookstore that will let users purchase best sellers from within the application
without having to open a browser.

\section{China Requires Censoring on New PCs}

BEIJING -- China has issued a sweeping directive requiring all personal computers sold in the
country to include sophisticated software that can filter out pornography and other ``unhealthy
information'' from the Internet.

The software, which manufacturers must install on all new PCs starting July 1, would allow the
government to regularly update computers with an ever-changing list of banned Web sites.

The rules, issued last month, ratchet up Internet restrictions that are already among the most
stringent in the world. China regularly blocks Web sites that discuss the Dalai Lama, the 1989
crackdown on Tiananmen Square protesters, and the Falun Gong, the banned spiritual movement.

But free-speech advocates say they fear the new software could make it even more difficult for
China's 300 million Internet users to obtain uncensored news and information.

``This is a very bad thing,'' said Charles Mok, chairman of the Hong Kong chapter of the Internet
Society, an international advisory group on Internet standards. ``It's like downloading spyware onto
your computer, but the government is the spy.''

Called Green Dam -- a reference to slogans that describe a smut-free Internet as ``green'' -- the
software is designed to filter out sexually explicit images and words, according to the company that
designed it. Computer experts, however, warn that once installed, the software could be directed to
block all manner of content or allow the government to monitor Internet use and collect personal
information.

Details of the new regulations, which were posted Monday on a government Web site, were first
reported by The Wall Street Journal.

PC makers that serve the Chinese market, among them Dell, Lenovo and Hewlett-Packard, said they were
studying the new rules and declined to comment. But privately, industry executives in the United
States said they were unnerved by the new rules, which were issued by the Ministry of Industry and
Information Technology with no consultation and no advance warning.

Beyond the nettlesome issue of abetting government censorship, they said six weeks was not enough
time to shift production on such a large scale. ``Many of us are going to take it in the neck with
this mandate,'' said one executive. ``It has put people into five-alarm mode.''

This is not the first time that foreign companies have been enlisted in government efforts to police
the Internet. Google already removes politically forbidden results yielded by its popular search
engine, Microsoft allows censors to block content on its blog service, and Yahoo was widely
criticized for turning over information that was used to jail a journalist.

``I would advise dissidents to buy computers before July 1,'' said Clothilde Le Coz, the head of the
Internet freedom desk of Reporters Without Borders.

More than 40 million personal computers were sold last year in China, one of the fastest growing
markets. Despite the slowing economy, industry analysts expect that figure to rise by 3 percent this
year.

A group of industry representatives met with American officials Monday to express their displeasure
with the new rules, said Susan N.~Stevenson, a spokeswoman for the United States Embassy in Beijing.
``We view any attempt to restrict the free flow of information with great concern,'' she said.

Zhang Chenming, general manager of Jinhui Computer System Engineering, a company that helped create
Green Dam, said worries that the software could be used to censor a broad range of content or
monitor Internet use were overblown. He insisted that the software, which neutralizes programs
designed to override China's so-called Great Firewall, could simply be deleted or temporarily turned
off by the user. ``A parent can still use this computer to go to porn,'' he said.

Although the directive is somewhat imprecise and suggests that manufacturers can provide the
software as a compact disc, it also says that it must be installed on computer hard drives as a
backup file. The five-point circular uses the word ``preinstall'' repeatedly and the first clause
unequivocally states: ``Imported computers shall preinstall the latest available version of the
'Green Dam' software before they are sold in China.''

Manufacturers complain that they have had been given little guidance by Chinese authorities. ``The
wording may be intentionally vague, but the message is clear: we have no choice in the matter,''
said one computer executive who spoke on condition of anonymity because some companies are hoping
they can persuade the government to ease the requirements.

Industry experts and civil libertarians say they are worried the software may simply be a Trojan
horse for greater Internet control. The software developers have ties to China's military and public
security agencies, they point out, and Green Dam's backers say the effort is supported by Li
Changchun, the country's chief propaganda official and a member of the decision-making body of the
Communist Party, the Politburo Standing Committee.

The software will be provided free, paid for by the government, and according to the official Green
Dam Web site, it has already been downloaded 3.2 million times. That figure includes thousands of
schools that were required to install the software by the end of May. The site claims that Chinese
manufacturers, including Lenovo, Inspur and Hedy, have already agreed to install 52 million copies
of the software on new computers.

In recent months China has tightened its Internet restrictions, including an ``antivulgarity''
campaign that has closed down thousands of pornographic sites but also nonsexual sites, including
some of the most popular bulletin boards and blog hosts. China already employs more than 30,000
censors and thousands who ``guide public opinion'' by flooding bulletin boards with comments
favorable to the Communist Party.

Last week, as the 20th anniversary of the military crackdown on the Tiananmen Square protests
approached, the government blocked a host of Internet services, including Twitter, Microsoft's
live.com and Flickr, a photo-sharing site, though by Monday evening, these sites had become
available again. YouTube has been inaccessible in China outside Hong Kong since March.

Even beyond ethical concerns, those who have tested the new software describe it as technically
flawed. An American software engineer said it led machines to crash frequently. Others worry that it
could leave tens of millions of computers vulnerable to hackers. So far, at least, there is no
version for the Linux operating system and for Apple's Macintosh system.

The directive makes no mention of the special administrative regions of Hong Kong and Macao, but one
industry association executive said companies had been told that these areas are exempt from the new
rules.

On Monday, Green Dam's own Web site offered a hint of discontent over the filtering software. On the
bulletin board section of the site, one writer described it as a ``Web devil'' and several users
complained that pornographic images slipped through or that their computers had become painfully
slow. ``It seems pretty lousy so far,'' one posting said. ``It's not very powerful; I can't surf the
Internet normally and it's affecting the operation of other software.''

By Monday night, however, most of the comments had been deleted.

\section{Opening Doors on the Way to a Personal Robot}

Consider it one small step -- or a roll, actually -- for a robot, one not giant, but significant
step for robotics.

Willow Garage, a Silicon Valley robotics research group, said that its experimental PR2 robot, which
has wheels and can travel at speeds up to a mile and a quarter per hour, was able to open and pass
through 10 doors and plug itself into 10 standard wall sockets in less than an hour. In a different
test, the same robot completed a marathon in the company's office, traveling 26.2 miles. PR2 will
not compete with humans yet; it took more than four days.

For the person who wants to buy a fully functioning robot butler, this may not seem so impressive.
But for roboticists and a new generation of technologists in Silicon Valley, this is a significant
achievement, a step along the way to the personal robot industry.

Willow Garage was founded by Scott Hassan, one of the designers of the original Google search
engine. The company's name is a reference to a small garage on Willow Road in Menlo Park, Calif.,
which was Google's first office. The company is trying to develop a new generation of robotic
personal assistants. Roboticists here and at other companies envision creating something on the
scale of the personal computer industry, with mechanical personal assistants taking over a lot of
drudgery, from cleaning up to fetching a beer from the refrigerator.

This is not a new hope, nor is it the first time that robots have tried to open doors, navigate
rooms and recharge themselves. The Beast, a robot built at Johns Hopkins University in the
mid-1960s, was able to locate standard wall sockets to refuel. And devices like the inexpensive
iRobot Roomba vacuum cleaner can locate and dock with a specially designed charging station.

But roboticists said that the Willow Garage robot was the first to integrate the ability to do a
number of operations in a real-world environment.

``There are other groups that have opened doors before,'' said Andrew Ng, a Stanford roboticist with
several students who have gone to work for the company. But, Mr.~Ng said, this seemed to be the
first robot able to repeatedly and reliably open doors and plug itself in.

William L.~Whittaker, a Carnegie Mellon University roboticist and the winner of a Defense Department
urban challenge robot driving contest last year, said it was ``unprecedented'' for a robot to
navigate in a building reliably and repeatedly recharge itself. ``These guys are the real deal,'' he
said.

The Willow Garage laboratory is less than a mile from another pioneering mobile robot project that
was developed at SRI International in the late 1960s. The robot, known as Shakey, had little onboard
computing capability and was remotely controlled by a combination of a mainframe and a minicomputer.

``In 40 years there has been a lot of progress, but not progress you notice,'' said Nils Nilsson, a
pioneering artificial intelligence researcher who was one of Shakey's designers. ``A lot of the
progress has been made in removing the cheats we used.''

To help Shakey navigate, he said, the baseboards in the room were painted black to help identify
walls, and objects were painted red so they could be identified by the vision system.

The ultimate goal of the Willow Garage researchers is to build a Robot Operating System, or R.O.S.,
that would greatly facilitate the work of a generation of software developers.

Microsoft is developing a similar system based on a version of the Windows operating system, but the
Willow Garage effort is an open-source project intended to leverage the contributions of a number of
robotics experts around the world.

Toward that end, a team of roboticists from the University of Tokyo recently modified the Willow
Garage R.O.S. to run on a robot they were developing, said Steve Cousins, Willow Garage's president
and chief executive.

``The eventual goal is to provide a set of capabilities that are so generic and so universal that
they can be used as building blocks in more complicated applications,'' said Sebastian Thrun,
director of the Stanford Artificial Intelligence Laboratory and a Willow Garage board member.

The current PR2 robot is powered by several Intel microprocessor chips and ``sees'' with a
combination of sensors including scanning lasers and video cameras. It is able to locate electrical
outlets and navigate in a building that is designed to accessibility guidelines in the Americans
With Disabilities act. Such buildings do not have doors with round door knobs, which simplifies the
task of opening a door for a robot. The PR2 carries its plug on a magnet at its base.

In the recent test, one of the plugs was behind a locked door, so the robot needed to be intelligent
enough to abort its effort and move on to the next plug. The doors it needed to navigate were
alternately closed, open and partially open.

One of the scientists summed up his feelings about the milestone in a succinct e-mail message:

``Now they can escape and fend for themselves.''

\section{Apple Races to Keep Ahead of Rivals}

SAN FRANCISCO -- With its coveted gadgets and resurgent stock price, Apple has cast something of a
spell on both consumers and investors.

At its annual Worldwide Developers Conference, which begins here on Monday, Apple executives will
try to sustain that magic, using what has become one of their highest-profile events of the year.

It will not be easy. While he is fond of surprises, Steven P.~Jobs, Apple's chief executive, whose
personal star power has been known to amplify the company's message, is not likely to be there. The
company says it continues to expect his return from medical leave at the end of the month.

Apple must also meet sky-high and perhaps unrealistic expectations from enthusiasts and analysts,
and it is just as likely to draw attention for what it does not unveil as for what it does. That
challenge is partly one of its own making. Apple's strategy -- particularly with the iPhone -- has
come to depend on a steady stream of hit devices that are viewed by consumers as being so far ahead
of the competition that they are worth paying extra money.

But the competition is now catching up. Palm, Google, Microsoft, Nokia and Research in Motion, maker
of the BlackBerry, are all at varying stages of developing and introducing their own iPhone-like
devices and software, along with easily accessible stores for the small programs known as
applications, or apps, that run on those devices. In some cases, those companies are releasing a
greater variety of phones, on more wireless carriers around the world, than Apple.

To maintain its advantage, Apple must preserve the impression that it is far ahead of rivals when it
comes to the capabilities and the ``cool'' factor of its devices.

``If they start making products people don't want, and start losing users, then Apple's strategy
will run into problems,'' said Benjamin Reitzes, an analyst at Barclays Capital. ``If they continue
to have an aura where their products are seen as defining the marketplace, they are going to be
fine. But that's going to be the challenge and the opportunity for Apple.''

For now at least, Apple appears to have a comfortable lead in the simmering smartphone battle. In
the last two years, it has sold more than 37 million iPhones and the similar-looking iPod Touches,
and the devices have become its most profitable product category. Apple's stock has nearly doubled
in the last six months, largely on the iPhone's inexorable momentum, although it is still down 28
percent from its high in December 2007.

The iPhone has also attracted an enthusiastic community of independent software makers that have
created about 40,000 free or low-cost applications for the devices. Apple now celebrates that wide
range of programs, from flight simulators to spreadsheets, as one of the major differences between
the iPhone and competing systems.

Apple's goal for next week: to keep the energy, momentum and spotlight focused on the iPhone and
away from its rivals, like the new Pre from Palm, a well-reviewed new phone that goes on sale
Saturday.

Apple is expected to use its keynote presentation on Monday to demonstrate the new types of programs
that can be created on the latest version of the iPhone's operating system.

Developers will be able to charge for certain features within their programs (perhaps a new level
within a game), and applications will be able to send alerts to users even whey they are not
actively running. Apple has also said it will show an early version of its latest operating system
for the Macintosh, dubbed Snow Leopard.

But talking about new software can be dull. So analysts point to Apple's need to surprise its fans,
and to its dwindling inventory of iPhones in stores, and predict that Apple will introduce a new
iPhone model next week as well. Its new tricks could include the ability to record video and an
internal compass that will add extra intelligence to the device's awareness of its location.

``Apple's objective will be to clearly and strongly show why the iPhone platform is the best for
developers, and that means there's a high probability Apple will announce one or more phones at the
conference,'' said Michael Abramsky, an analyst at RBC Capital Markets.

Satisfying its large audience of developers may be Apple's most significant challenge. They want to
know they can build profitable businesses even if they focus exclusively on the iPhone, and that
Apple will continue to do its part by wooing consumers with compelling new gadgets.

``We want to hear that there are going to be cool new devices, because that means our market gets
bigger,'' said Bart Decrem, the chief executive of Tapulous, a company in Palo Alto, Calif., that
makes musical games for the iPhone.

Apple has actually set a high bar for itself. Unlike its rivals Google and Microsoft, which license
their mobile operating systems to many phone makers, Apple builds its own hardware and software and
carefully strikes exclusive relationships with wireless carriers that are willing to heavily
subsidize its devices. The strategy depends on a constant flow of new products that people are
willing to switch wireless companies and pay extra to use.

That same strategy -- introducing expensive but elegant products with high switching costs -- is
showing signs of strain in other parts of Apple's business. Although Apple has performed better
lately than other American computer makers, its sales of Mac laptops and desktops declined by 3
percent in the first three months of the year. That was the first time in six years that sales in
Apple's personal computer business had a year-over-year decline.

Meanwhile, the Taiwanese manufacturers of the smaller, cheaper computers called netbooks, like Acer
and Asustek, continue to have growing sales.

``Apple has to address the structural slide of the PC industry to more value-priced products,'' said
Ashok Kumar, an analyst at Collins Stewart.

But Apple is not expected to talk about the Mac next week, and that could be emblematic of one of
the largest problems the company faces. Apple watchers and investors have many questions that Apple
does not seem to want to answer in public, and they could punish the stock if they come away from
the conference disappointed.

The big questions include whether Apple is working on any entirely new product lines (like a much
rumored tablet computer), what it might do with its \$29 billion hoard of cash, and how Mr.~Jobs's
health is faring.

Mr.~Jobs, Apple's innovator in chief, has not been seen in public since last October, when he
unveiled a line of laptops at Apple's headquarters. A survivor of pancreatic cancer, Mr.~Jobs has
been on medical leave since January. Two people close to his personal circle say he looks healthier,
and they expect him to be back at work full time soon. Apple will say only that it still expects
Mr.~Jobs back at the end of June.

That is one area where the expectations of Apple observers and investors may have fallen somewhat.
Apple has thrived in the last six months under the leadership of Timothy D.~Cook, Apple's chief
operating officer, who has handled day-to-day operations during Mr.~Jobs's absence. If Mr.~Jobs does
not return, or returns in a diminished role, investors and fans may no longer see it as the end of
the world.

``Apple has had a nice rally because they put up very strong numbers, and at the end of the day it
will still be all about numbers,'' said Mr.~Reitzes of Barclays Capital.

\section{Chinese Mayor Detained in Graft Inquiry}

BEIJING -- The mayor of Shenzhen, one of China's most dynamic cities, has been detained as part of a
graft investigation that threatens to overturn the local political establishment, the state media
reported Tuesday.

The mayor, Xu Zhongzheng, who is also deputy secretary of Shenzhen's Communist Party, failed to show
up for a public appearance on Friday. On Monday, officials said that he was being held under a form
of detention reserved for party members.

Although it did not provide much detail on the case, Xinhua, China's official news agency, reported
that Mr.~Xu was being questioned about ``serious disciplinary offenses,'' which is often a prelude
to criminal charges. A magazine, Caijing, said Mr.~Xu's wife was also being held.

The Communist Party has been struggling to tamp down corruption, although it is not always clear
whether those caught up in antigraft campaigns are also victims of high-level political
machinations. Last year, the former Communist Party chief of Shanghai, Chen Liangyu, was sentenced
to 18 years in prison for bribery and fraud, but many here viewed his downfall through the prism of
a power struggle involving China's top leaders.

A onetime car technician who steadily rose through the party, Mr.~Xu, 54, has been an aggressive
booster for Shenzhen, a fast-growing city of eight million near Hong Kong that was little more than
a fishing village when it became China's first special economic zone in 1980. According to reports
in the Hong Kong media, Mr.~Xu may have been singled out for his connections to Huang Guangyu, a
appliance store magnate and billionaire who is being held on charges of stock manipulation and
``economic crimes.''

Mr.~Huang's arrest last November has led to more investigations that have ensnared a number of
senior officials. One of them, Zheng Shaodong, the assistant minister for anticorruption efforts in
Guangdong Province, was removed from his job for his role in the scandal, Xinhua reported Tuesday.

After he was appointed mayor in 2005, Mr.~Xu pushed for bureaucratic changes and reductions in
government spending, but he was criticized by some for failing to intervene in the city's overheated
real estate market, which has since crashed.

Xinhua on Tuesday described him as a ``hard-liner'' and implied that his detention might have a
political tinge. ``His blunt refusal to put a ceiling on the local housing market caused nationwide
controversy over whether he was representing the interests of major real estate developers instead
of the public, who found housing prices unaffordable,'' the news agency said.

Allegations of public corruption are nothing new in Shenzhen. In 2006, two dozen local officials
were detained as part of a bribery scandal. Among them was the vice president of Shenzhen's
Intermediate People's Court, who was accused of leading a ``decadent lifestyle'' and of trading
favorable verdicts for money.

Last year, Yu Youjun, a former mayor of Shenzhen who later became a top official in the Ministry of
Culture, was passed over for promotion and removed from the Communist Party's Central Committee after a disciplinary procedure.

\section{Smartphone Rises Fast From Gadget to Necessity}

In today's recession-racked economy, penny-pinching is a national pastime. But people are still
opening their wallets for smartphones.

Sales of BlackBerrys, iPhones and other smartphone models are rising smartly and are projected to
increase 25 percent this year, according to Gartner, a research business. Widely anticipated new
models like the Palm Pre, which went on sale nationwide on Saturday, will help fuel that growth.
Meanwhile, total cellphone sales are expected to fall.

The smartphone surge, it seems, is a case of a trading-up trend in technology that is running strong
enough to weather the downturn. And as is so often true when it comes to adoption of new technology,
the smartphone story is as much about consumer sociology and psychology as it is about chips, bytes
and bandwidth.

For a growing swath of the population, the social expectation is that one is nearly always connected
and reachable almost instantly via e-mail. The smartphone, analysts say, is the instrument of that
connectedness -- and thus worth the cost, both as a communications tool and as a status symbol.

``The social norm is that you should respond within a couple of hours, if not immediately,'' said
David E.~Meyer, a professor of psychology at the University of Michigan. ``If you don't, it is
assumed you are out to lunch mentally, out of it socially, or don't like the person who sent the
e-mail.''

The spread of those social assumptions may signal a technological crossover that echoes the
proliferation of e-mail itself more than a decade ago. At some point in the early 1990s, it became
socially unacceptable -- at least for many people -- to not have an e-mail address.

Smartphones are not cheap, particularly in tough economic times. The phones, even with routine
discounts from wireless carriers, usually cost \$100 to \$300, while the data and calling service
plans are typically \$80 to \$100 a month.

But recent smartphone converts are often people who count pennies, including many from the growing
ranks of job seekers. Helene Rude of Briarcliff Manor, N.Y., was laid off from her job as a business
development manager at I.B.M. this year, when her unit, among others, was the target of cuts. When
she left, Ms.~Rude had to turn in her company notebook computer with its constant wireless
connection.

So she got an iPhone instead, allowing her to be online no matter where she was, without having to
lug a computer around. ``I absolutely got it for the job search,'' she said. ``I don't know if it's
really an expectation, but if another job candidate returns an e-mail message eight hours later, and
you get back immediately with a message that says 'Sent from my iPhone,' I think it has to be a
check box in your favor.''

That is certainly the sort of message the wireless industry would like to reinforce. ``Smartphones
are seen as essential to be productive in a mobile society,'' said David Christopher, chief
marketing officer at AT\&T's wireless division.

James L.~Balsillie, co-chief executive of Research in Motion, maker of the BlackBerry, said the
company's introduction of less business-oriented phones, with the general spread of mobile
communication, explained the snowballing growth in BlackBerry users. They now number 25 million,
nearly double the total a year ago and a tenfold increase in the last four years.

The smartphone wave, industry analysts say, should continue to build. The room for gains is ample
because, though rising, smartphone sales will still account for only a quarter of total cellphone
shipments in the United States this year. And along with the Palm Pre, a host of new smartphone
handset and software offerings are coming this year, from Apple, R.I.M., Nokia, Microsoft, Google
and others.

The industry's goal is to win over more rank-and-file converts like Joseph Sexton of San Jose,
Calif., who calls himself ``not a gadget person.'' Mr.~Sexton, 45, decided to leave his job as a
manager of a community health organization to travel and then look for other work, shortly before
the financial crisis hit last fall.

He soon found himself looking for work in a tough market, and got a smartphone as a digital
assistant in his search. Now he is hooked. ``It allows me to be on top of things, and always
connected, no matter where I am,'' he said.

Mr.~Sexton searches the Web, takes notes and sends e-mail with his iPhone. When he has trouble
sleeping, he reaches for his smartphone to read news or check e-mail. In fact, Mr.~Sexton said, he
finds himself reading more online these days and buying fewer magazines.

``Basically, I'm walking around with a minicomputer in my pocket,'' he said. ``And it's a part of me
now, an appendage.''

Such a digital connection can have its downside. The perils of obsessive smartphone use have been
well documented, including distracted driving and the stress of multitasking. CrackBerry, a term
coined years ago, is telling.

The smartphone, said Mr.~Meyer, a cognitive psychologist, can be seen as a digital ``Skinner box,''
a reference to the experiments of the behavioral psychologist B.~F. Skinner in which rats were
conditioned to press a lever repeatedly to get food pellets.

With the smartphone, he said, the stimuli are information feeds. ``It can be powerfully reinforcing
behavior,'' he said. ``But the key is to make sure this technology helps you carry out the tasks of
daily life instead of interfering with them. It's about balance and managing things.''

\end{document}
